\documentclass[../main.tex]{subfiles}

\begin{document}

\chapter{ Week 1 }

代靖涵 25120222201319

\begin{exercise}{}{}
A nonassociative algebra $(A, \circ)$ is called \dfntxt{left symmetric} if the associators are left symmetric in the sense that
\[\begin{aligned}
(x,y,z) = (y,x,z) \quad ((x,y,z) := (x \circ y) \circ z - x \circ (y \circ z))
\end{aligned}\]
for all $x,y,z \in A$. Prove that for every left symmetric algebra $(A, \circ)$, $(A, [\cdot, \cdot])$ is a Lie algebra, where
\[\begin{aligned}
[\cdot, \cdot] : A \times A \to A, \quad (x,y) \mapsto [x,y] := x \circ y - y \circ x
\end{aligned}\]
denotes the commutator on $A$.
\end{exercise}

\begin{proof}
    To show \(  \left[ \cdot ,\cdot  \right]   \) is bilinear, we  see  \[
    \begin{aligned}
    \left[ x+ y,z \right]&= \left( x+ y \right)\circ z- z\circ\left( x+ y \right)    \\ 
     &= \left( x\circ z - z\circ x \right) +\left(  y\circ z - z \circ y \right) \\ 
      &=  \left[ x,y \right]+ \left[ y,z \right]  
    \end{aligned}
    \]  Similarly \[
    \begin{aligned}
    \left[ x,y+ z \right]&= x\circ \left( y+ z \right)- \left( y+ z \right)\circ x\\ 
     &= \left( x\circ y- y\circ x \right)+ \left( x\circ z-z\circ x \right)\\ 
      &= \left[ x,y \right]+ \left[ y,z \right]        
    \end{aligned}
    \]

    Futhermore, it is obvious that \[
    \left[ x,x \right]= x\circ x-x\circ x= 0 
    \]
  Then we only need to show the Jacobi identity. 
 Since \(  \left( A,\circ  \right)   \) is left symmetric, we have 
    \[
    \begin{aligned}
    & \left[ x\left[ yz \right]  \right]+ \left[ y\left[ zx \right]  \right]+ \left[ z\left[ xy \right]  \right]   \\ 
     &= x\circ \left( y\circ z - z\circ y \right)- \left( y\circ z - z \circ y \right)\circ x\\ 
      &+  y\circ \left( z\circ x - x \circ z \right)-\left( z\circ x - x \circ z  \right)\circ y\\ 
       &+  z\circ \left( x\circ y-y\circ x \right)-\left( x\circ y-y\circ x \right)\circ z\\ 
        &= x\circ \left( y\circ z \right)- x\circ \left( z\circ y \right)-\left( y\circ z  \right)\circ x + \left( z\circ y \right)\circ x\\ 
         &+ y\circ \left( z\circ x \right)-y\circ \left( x\circ z \right)-\left( z\circ x  \right)\circ y             + \left( x\circ z \right)\circ y\\ 
          &+ z\circ \left( x\circ y \right) -z\circ \left( y\circ x \right)-\left( x\circ y \right)\circ z + \left( y\circ x \right)\circ z\\ 
           &=  \left( x\circ \left( y\circ z \right)-\left( x\circ y \right)\circ z   \right)- \left( x\circ \left( z\circ y \right)+ \left( x\circ z \right)\circ y   \right)\\ 
            &- \left( \left( y\circ z \right)\circ x  + y\circ \left( z\circ x \right) \right)+ \left( \left( z\circ y \right)\circ x-z\circ \left( y\circ x \right)   \right)\\ 
             &- \left( y\circ \left( x\circ z \right)- \left( y\circ x \right)\circ z   \right)-\left( \left( z\circ x \right)\circ y -z\circ \left( x\circ y \right)   \right)\\ 
              &= 0       
    \end{aligned}
    \] Thus \(  \left( A,\left[ \cdot ,\cdot  \right]  \right)    \) is a Lie algebra. 
\end{proof}

\begin{exercise}{}{}
Write $\operatorname{Der}(A)$ for the Lie algebra of derivations on $A$, where $A = \mathbb{C}[t, t^{-1}]$. Prove that $\operatorname{Der}(A)$ has a basis $\{L_n := t^{n+1} \frac{d}{dt} \mid n \in \mathbb{Z}\}$ and the Lie brackets among the basis elements are as follows:
\[\begin{aligned}
[L_n, L_m] = (m-n)L_{m+n} \quad (m,n \in \mathbb{Z}).
\end{aligned}\]
\end{exercise}

\begin{proof}
 
     Let \(  D\in \mathrm{Der}\left( A \right)   \) be any derivation on \(  A  \). Then \[
    \begin{aligned}
     D\left( 1  \right)&= D\left( t\cdot t^{-1}  \right) \\ 
      &=  D\left( t \right)\cdot t^{-1} + tD\left( t^{-1}  \right)   
    \end{aligned}  
     \]  The other hand \[
     D\left( 1 \right)= D\left( 1\cdot 1 \right)=  D\left( 1 \right)\cdot 1+ 1\cdot D\left( 1 \right)= 2D\left( 1 \right)   \implies D\left( 1 \right)= 0  
     \] Thus we have \[
     D\left( t \right)\cdot  t^{-1} + tD\left( t^{-1}  \right)= 0  \implies D\left( t^{-1}  \right)=  -t ^{-2}D\left( t \right) 
     \] From the linearity and the product rule, we observe that a dirivation \(  D \in \mathrm{Der}\left( A \right)   \) is totally determined by its value on  \(  t  \), that is  for any \(  D_1,D_2 \in \mathrm{Der}\left( A \right)   \) if \(  D_1\left( t \right)= D_2\left( t \right)    \), then \(  D_1= D_2  \).    It is straightforward to verify that  \(  L_{n}  \) is a derivation on \(  A  \) for each \(  n \in \mathbb{Z}   \). Note that \[
     L_{n}\left( t \right)=  t^{n+ 1} ,\quad  n\in \mathbb{Z} .
     \]   Since \(  D\left( t \right) \in A   \), we have  \[
     D\left( t \right)\in \operatorname{span}\left(\cdots ,L_{-2}\left( t \right),L_{-1}\left( t \right),L_0\left( t \right),L_1\left( t \right),L_2\left( t \right),\cdots       \right)  .
     \] Then we get \[
     D\in \operatorname{span}\left( \cdots ,L_{-2},L_{-1},L_{0},L_{1},L_{2},\cdots  \right) .
     \] Suppose we have a finite linear combination of the $L_n$ that equals the zero derivation:
$\sum_{n \in I} k_n L_n = 0,$
where $I \subset \mathbb{Z}$ is a finite set and $k_n \in \mathbb{C}$.
To check whether this is the zero derivation, we can test its action on the generator $t$.
$\left(\sum_{n \in I} k_n L_n\right)(t) = 0(t) = 0.$
This gives:
$\sum_{n \in I} k_n L_n(t) = \sum_{n \in I} k_n t^{n+1} = 0.$
The set of monomials $\{t^j\}_{j \in \mathbb{Z}}$ is a basis for $A$ as a $\mathbb{C}$-vector space, so they are linearly independent. The equation $\sum_{n \in I} k_n t^{n+1} = 0$ implies that all coefficients must be zero, i.e., $k_n = 0$ for all $n \in I$.
This proves that the set $\{L_n\}_{n \in \mathbb{Z}}$ is linearly independent. Combining above all , we know \(  \mathrm{Der}\left( A \right)   \) has a basis \(  \left\{ L_{n}:n\in \mathbb{Z}  \right\}  \).      

     To show \(  \left[ L_{n},L_{m} \right]= \left( m-n \right)L_{m+ n}    \), we only need to show that they agree on \(  t  \): \[
     \begin{aligned}
      \left[ L_{n},L_{m} \right]\left( t \right)&=  \left( L_{n} \left( L_{m}\left( t \right)  \right)-L_{m}\left( L_{n}\left( t \right)  \right)   \right)\\ 
       &=    \left( L_{n}\left( t^{m+ 1} \right)-L_{m}\left( t^{n+ 1} \right)   \right)\\ 
        &= \left( \left( m+ 1 \right)t^{m+ n+ 1}  -  \left( n+ 1 \right)t^{m+ n+ 1} \right)\\ 
         &= \left( m-n \right) t^{m+ n+ 1}\\ 
          &= \left( m-n \right)L_{m+ n}\left( t \right)     
     \end{aligned}
     \]   which completes the proof.

    
\end{proof}

\begin{exercise}{}{}
Let $\mathbb{F} = \mathbb{R}$ or $\mathbb{C}$, and let $V$ be an \dfntxt{n-dimensional} vector space over $\mathbb{F}$ equipped with a nondegenerate bilinear form
\[\begin{aligned}
B: V \times V \to \mathbb{F}, \quad (v,w) \mapsto B(v,w).
\end{aligned}\]
Prove that:
\begin{enumerate}
\item[(a)] If $\mathbb{F} = \mathbb{C}$ and $B$ is symmetric, then there exists a basis $\{v_1, v_2, \dots, v_n\}$ of $V$ such that $B(v_i, v_j) = \delta_{i,j}$ for all $1 \le i, j \le n$.
\item[(b)] If $\mathbb{F} = \mathbb{R}$ and $B$ is symmetric, then there exists a basis $\{v_1, v_2, \dots, v_n\}$ of $V$ and $p=0, 1, \dots, n$ such that $B(v_i, v_j) = \varepsilon_i \delta_{i,j}$ for all $1 \le i, j \le n$, where $\varepsilon_i = 1$ if $i \le p$ and $\varepsilon_i = -1$ if $i > p$.
\item[(c)] If $\mathbb{F} = \mathbb{R}$ or $\mathbb{C}$ and $B$ is skew-symmetric, then $n=2m$ is even and there exists a basis $\{v_1, v_2, \dots, v_n\}$ of $V$ such that the matrix
\[\begin{aligned}
(B(v_i, v_j))_{1 \le i, j \le n} = \begin{pmatrix} 0 & I_m \\ -I_m & 0 \end{pmatrix},
\end{aligned}\]
where $I_m$ denotes the identity matrix of rank $m$.
\end{enumerate}
\end{exercise}

\begin{proof}
    \item (a) 
\end{proof}
\end{document}